% ============================================================
\section{Analysis and Discussion}
\label{sec:analysis}
% ============================================================

\textbf{Why does base Moirai fail on \dataset{}?}
The zero-shot Moirai model (condition A, ROC-AUC $\approx$ 0.27) fails
because it was pre-trained on financial, energy, and weather time series.
IoT network traffic has fundamentally different statistical properties:
high autocorrelation in packet inter-arrival times, bursty protocol
behaviour, and feature scales that differ from Moirai's pre-training
distribution.  Fine-tuning on just 10 epochs of domain-specific data
substantially closes the distribution gap (condition D vs A).

\textbf{Why do hard negatives help more than real attacks?}
Standard attack traffic in \dataset{} is high-magnitude (large Z-scores,
easily detected by simple thresholds).  These easy negatives let the
model converge to a coarse decision boundary.  Synthetic hard negatives
at stealth 85--95\% are deliberately positioned near the boundary,
forcing the model to learn fine-grained temporal features.
This is analogous to curriculum learning~\cite{bengio2009curriculum}
and hard example mining in object detection~\cite{shrivastava2016training}.

\textbf{Protocol constraint utility.}
Beyond improving detection quality, the constraint validation layer
provides an interpretable quality metric for generated samples.
The constraint pass-rate statistics (Table~\ref{tab:constraints})
could serve as an independent signal for deployment engineers to
verify that synthetic training data reflects real IoT traffic.

\textbf{Computational overhead.}
Hard-negative generation with Diffusion-TS (mock mode) takes $<$2s
per sample.  Fine-tuning Moirai-Small for 10 epochs on our dataset
requires $\approx$2--4 GPU-hours on an NVIDIA A100.
Inference latency is $<$50ms per sample, suitable for near-real-time
IoT monitoring at moderate sampling rates.

\textbf{Limitations and future work.}
(1) Our evaluation uses a single dataset; generalisation to other
IoT benchmarks (e.g., N-BaIoT, UNSW-NB15) is future work.
(2) Moirai requires a fixed 12-dimensional input; adapting to
variable-feature IoT environments would require input projection.
(3) The diffusion model operates in mock mode (statistical generation)
when Diffusion-TS is not installed; full quality requires the full
model.  (4) The stealth thresholds (85--95\%) are fixed; adaptive
adversaries could target higher stealth levels.

\textbf{Responsible disclosure.}
The same diffusion-based adversarial augmentation technique used here
for defensive training could potentially be used offensively to evade
existing IDS deployments.  We follow responsible disclosure norms:
the code and synthetic data are released to enable defensive research
and reproducible benchmarking, not to enable attacks.
